\documentclass[11pt]{article}
\usepackage[margin=1in]{geometry}
\usepackage{amsmath,amssymb,booktabs,enumitem}
\newcommand{\fn}[1]{\texttt{\detokenize{#1}}}
\title{Step 3, v3.2: the far-field localisation gap --- precise statement and feasibility of two repairs}
\author{}\date{12 November 2026}
\begin{document}\maketitle
\noindent\textbf{Status.} Feasibility pass only. Nothing here is a proof; numbers marked ``est.'' are float estimates, the others are certified elsewhere as cited. The repair was \emph{not} carried out: both routes need an ingredient that is not established (\S\ref{sec:stop}).

\section{The gap (for the main-theorem document)}\label{sec:gap}
\noindent\fbox{\parbox{0.97\linewidth}{\textbf{Gap statement.} Step 3 controls the far field $y\ge60$ (rescaled) by energies on windows fixed in physical space (Lemma~S(B)). Its certified inequalities (\fn{lemmaS_B.py}, \fn{lemmaS_F4B.py}, \fn{lemmaS_Bbox.py}: ``no ramp inflow'') contain the transport terms and the interior fourth-order terms, but \emph{not the fourth-order terms at the edges of a window}, and the outer end of the shell energies of Lemma~S(A) is likewise unquantified (``outer cutoff far away''; no outer-end term in \fn{lemmaS_F1.assemble}, where \fn{extra_outer} is the inner-ramp foot/S33 term). Estimating an edge by a cutoff requires, on the cutoff ramp and for all later times, bounds on $z_k=(y\partial_y)^kv$ for $k\le j+1$ at level $j$; at level 2 this is a space--time bound on $z_3$ (four derivatives of the curve). v3.2 supplies it at the outer edge, inside the Lemma~O region (Lemma~R, certified). It is \textbf{unproved at every edge lying in $60\lambda(\tau)\lesssim r<150\lambda_1$}, at levels 1 and 2 (level 0 needs only $z_1$). Affected: (i) the strip statement, i.e.\ hand-off-level energies on $r\in[50,150)\lambda_1$ for all $t$ (levels 1, 2), hence Lemma~O hypothesis H6 and the bootstrap arrows D4, D5; (ii) the crude window bounds beyond $y=60$ ($E_B$, $|f|\le\frac12$, curvature $\le S_1^f/r$: B3, W4), which v3.1 derives by the same edge-free computation; (iii) through (ii), the coefficient box and the outer end of Lemma~S(A), which v3.1 uses in the form ``$|V|\le\frac12$ beyond $y=60$''. Not affected: the linear theory (steps 1--2 and the coercivity constants), Lemma~O and Lemma~R as local lemmas, the verifier \emph{given} its inputs $S_0^v,S_1^v$ and the flux constants, and the core and inner-ramp computations of Lemma~S(A) on $y\in[14,60]$ \emph{given} a crude box ($|f|\le0.11$, $y|f'|\le0.52$) and an $O(1)$ space--time bound for $z_2,z_3$ on its outer cutoff. So the main trapping argument is independent of the \emph{hand-off-level} strip bounds and of H6, but it is \textbf{not} independent of the gap: it needs the crude far-field bounds (ii), with large numerical room (core margins $5.9$ and $2.8$) but without a complete proof.}}

\section{Why a wide ramp alone does not work (route Q)}
A cutoff ramp contributes to the level-$j$ energy at most $\sum_{k\le j+1}c_{jk}\bar N_k^2$, with $\bar N_k^2$ the average over physical time of $\int_{\rm ramp}z_k^2ds$ (frozen normalisation). Certified coefficients (\fn{lemmaS_v32/ramp_in_50_75.json}, \fn{ramp_out_*.json}): inner ramp $[50,75]$: $c_{22}=2.4\cdot10^{-4}$, $c_{23}=3.4\cdot10^{-5}$; they scale like $y_0^{-5}\ell^{-2}$ to $\ell^{-4}$ ($\ell$ the log-width), so a ramp $[50,150]$ gains a factor $7$--$50$.
\begin{itemize}[leftmargin=*,itemsep=1pt]
\item The obstacle is not the size of $c_{jk}$ but $\bar N_3$: inside $r<150\lambda_1$ there is \emph{no} time-uniform source for $z_3$ on a ramp fixed in physical space. Lemma~R needs cylinders with a-priori bounds and time ratio below $t^*$; the dissipation of Lemma~S(A) controls $\varepsilon\omega z_3^2$ only at fixed rescaled $y$, and for a physical ramp it degrades like $e^{(\tau-\tau_1)/4}$ after weighting with $dt$ (est.).
\item With the cone geometry of v3.1/v3.2 ($L=r/4$) the collar that needs control always starts at $r=L_c$, where the tolerance is small. Fixed $L=50$, $X=3$, centres $\ge300\lambda_1$ (cylinders $[150,450]$, conclusion regions from $250$; H4 error part $\approx10^{-3}M^2$, est.\ from \fn{h4loc}; all unit ratios $1$, far family $113/50$) moves the collar to $[150,250)$, where the conversion $(50/y)^3$ raises the tolerated unweighted level-2 window energy to $\approx0.3$ (est.), against $8\cdot10^{-5}$ at hand-off. So Q makes \emph{crude} control of the collar sufficient, but does not produce it.
\end{itemize}

\section{Route P: weights that never vanish}
Replace the inner cutoff by a co-moving power weight $\Theta_2=(r/Y_1)^\mu$, $Y_1=150\lambda_1$, capped at $1$ for $r\ge Y_1$, on the frozen-profile deviation $v=f-U_1$ (Prop.~B$'$: forcing negligible). Then no $z_3$ is needed from outside: the weight's derivatives are bounded by the weight, and the only cutoff is a ramp fixed in rescaled coordinates inside the core of Lemma~S(A), $y\in[28,56]$, where (A) supplies $z_{0,1,2}$ and, through its level-1 dissipation, $z_3$. The non-integrable rate $\varepsilon(y)\hat G$ is removed by adding $g$ times the energy with exponent $\mu-4$ (since $\varepsilon y^\mu=y^{\mu-4}$, and that energy is damped at rate $1$). Both are pure-power weights, so the divergence-form tables of \fn{lemmaS_F1} apply with $\gamma=-\mu$ and $\gamma=4-\mu$.

\emph{Float numbers} (\fn{feas_P.py}, box $|f|\le0.11$, $y|f'|\le0.52$, no $L^4$ extras): for $\mu=4$: $\hat G_{jj}=(21,\,302,\,2.7\cdot10^3)$, so $\varepsilon(28)\hat G_{22}=0.0044$; ramp constants relative to $\varepsilon\Theta$: $K_{2\cdot}=(0,\,582,\,3.1\cdot10^3,\,704)$. Budget for the level-2 energy, unweighted on $r\ge150\lambda_1$ (est.): hand-off data $0.009$ (dominated by the shell threshold $E_2\le0.2$ on $[28,50]$ times the weight $(50/150)^4$), inflow through the ramp $0.013$, ramp commutators $0.004$ ($z_3$ from the level-1 dissipation of (A): $\le1.3\cdot10^{-4}$ per unit $\tau$ in the $\varepsilon\omega$ norm). Total $\approx0.03$ against a tolerance $\approx0.3$; then collar $K_2\approx0.03$, $K_3\approx0.03$ in units $L=50$ against $0.18$.

\section{What stops the repair}\label{sec:stop}
P combined with the geometry of Q closes on float estimates with a margin of about $10$. It was not proved, for two reasons.
\begin{enumerate}[leftmargin=*,itemsep=1pt]
\item \textbf{Crude box on the support of the weight.} The rates above need $|f|\le0.11$, $y|f'|\le0.52$ on the whole support, i.e.\ also on $60\lambda(\tau)\le r\le50\lambda_1$ at late times. The energy with $Y_1=150\lambda_1$ gives there only $y|f'|\lesssim1.4$ at $r=25\lambda_1$ (est.), outside the box; energies launched at later times and smaller scales start from the crude level of (A) and, with a one-sided power weight, lose a factor growing with the number of windows outside them. A two-sided (uniformly local) weight is needed, which is not a sum of powers, so the existing tables do not apply. This is item (ii) of the gap statement, and it is the missing ingredient for \emph{both} routes.
\item \textbf{Scope.} Even granting 1, Lemma~P needs new certified assemblies (two weights, the ramp $[28,56]$ fed by (A), the $L^4$ classes), a new innermost Lemma~O family, and step-4 reruns for it.
\end{enumerate}
\textbf{Proposed next step.} A uniformly local version of Lemma~S(B): general weights $W$ with $|W^{(m)}|\le\Gamma_mW$ in the divergence-form tables (the weight enters them linearly, so $\gamma^m$ becomes $\Gamma_m$ with absolute values), launched at every scale, giving first the crude box (ii) and then, with $Y_1=150\lambda_1$ and the fixed-$L$ family, the collar bounds for H6.

\section{Files}
\fn{feas_P.py}, \fn{feas_P.json} (this directory; run from \fn{writeup/lemmaS_v32}); certified inputs: \fn{lemmaO_v32/lemmaR.tex}, \fn{lemmaS_v32/ramp_cert.json}, \fn{ramp_in_50_75.json}.
\end{document}
